\chapter{Une interface web en localhost}

Un projet de recherche finit par manier plus de documents qu'on n'en peut ouvrir à la
main : des scans, des transcriptions, des tableaux de résultats, des listes de questions.
Plutôt que d'empiler les fenêtres, on peut demander à l'assistant de construire une petite
application web qui tourne sur sa propre machine, à l'adresse \texttt{localhost}. Elle n'a
pas besoin d'être belle ni d'être publique ; elle doit montrer côte à côte ce qu'on veut
comparer, et laisser une trace de ce qu'on a décidé.

\section{L'outillage}

On n'a pas besoin d'être développeur web. L'assistant installe un \emph{dev server}
(Vite, avec React, dans le projet qui sert d'exemple), écrit les pages, et le lance :

\begin{terminal}
npm install
# lance le serveur ; la page se recharge à chaque modification
npm run dev
# puis ouvrir http://localhost:5173
\end{terminal}

L'application de bureau de Claude a un navigateur intégré. L'assistant y ouvre lui-même la
page, lit la console, fait une capture d'écran et vérifie qu'un bouton fait ce qu'il doit
faire avant de dire que c'est fini. On travaille donc en conversation : « mets le scan à
gauche et la transcription à droite », « coche ce filtre par défaut », et l'on voit le
résultat dans la seconde.

\begin{regle}
Une page locale lit des fichiers du dépôt ou des fichiers ignorés par \fichier{.gitignore}
; elle n'envoie rien nulle part. Ce qui est privé le reste, même quand la page l'affiche.
\end{regle}

\section{Consulter plusieurs fonds à la fois}

Le premier usage est la lecture. Une page par document, avec à gauche l'image de la page
servie par l'archive qui la détient, et à droite la transcription, la lecture moderne, les
notes. Un onglet par fonds, un catalogue filtrable, un lien direct vers chaque page.

\begin{exemple}
Dans le projet, les images des archives ne sont jamais copiées sur le site public : la page
les demande à l'archive au moment de la lecture, par un petit relais, et affiche « ne
répond pas » quand l'archive est en panne. En local, les copies privées servent à
transcrire, jamais à être affichées à d'autres.
\end{exemple}

Un interrupteur décide de ce que la page montre. Les transcriptions d'un fonds dont on n'a
pas l'autorisation s'affichent en local et disparaissent du site public par un seul réglage
(chapitre~\ref{ch:droits}).

\section{Faire des reviews}

Le deuxième usage est la vérification humaine (chapitre~\ref{ch:review}). L'assistant
dépose ses doutes dans un fichier JSON ; la page de review en fait une liste de questions
fermées, chacune avec l'image de la page et la lecture proposée. On répond par des
boutons, on ajoute une note si besoin, et un bouton « copier les réponses » produit un bloc
de texte qu'on colle dans la conversation. L'assistant applique les réponses, retire les
questions réglées et reformule celles qui restent.

\begin{tipblock}
Les réponses restent dans le navigateur (\texttt{localStorage}) tant qu'on ne les a pas
envoyées. On peut fermer la page au milieu d'une séance et la reprendre le lendemain.
\end{tipblock}

La même page sert à préparer une visite en bibliothèque : une fiche par boîte d'archives,
les questions à poser à chaque pochette, et des cases à cocher qu'on remplit sur place,
depuis un téléphone.

\section{Construire des simulations}

Le troisième usage est l'exploration. Quand une question se teste sur des exemples, on
demande à l'assistant une page interactive : des curseurs pour les paramètres, une figure
qui se redessine, un tableau de résultats. Le calcul lourd reste dans un script (Python,
C, Lean) dont la page lit la sortie ; la page sert à regarder, à repérer le cas qui
cloche, à formuler la conjecture suivante.

\begin{exemple}
Pour un problème d'arrangements de droites, un script énumère tous les cas jusqu'à neuf
droites et écrit ses résultats dans un fichier. Une page peut alors montrer chaque
arrangement, avec ses sommets et ses triangles, et colorer ceux où une hypothèse tombe en
défaut. Regarder ces cas apprend souvent plus que lire les chiffres : on voit quelle
hypothèse supplémentaire manque.
\end{exemple}

\section{Ce qu'on garde de ces pages}

Une page locale est un instrument de travail, pas un produit. On la garde simple, on la
modifie en conversation, et l'on n'hésite pas à la jeter quand la question est réglée.
Ce qui doit durer passe ailleurs : les réponses dans les fichiers du dépôt, les résultats
dans un commit, les conclusions dans une issue ou un article. Quand une page mérite d'être
partagée, elle part sur le site public, derrière les mêmes interrupteurs.
